% !TeX program = xelatex
\documentclass[11pt,twoside]{book} % twoside = opposing margins
% --- METADATA (put in preamble) ---
\newcommand{\BookTitle}{An Easter Disciple:}
\newcommand{\BookSubtitle}{The Chronicle of Quintus, the Roman Knight}
\newcommand{\BookAuthor}{Arthur Benton Sanford}
\newcommand{\PublisherName}{Fly By Night Publishing}
\newcommand{\PublisherCity}{Queen Creek, AZ}
\newcommand{\EditionNote}{Text sourced from Project Gutenberg. This edition has been reformatted and lightly edited for clarity and modern typographic conventions.}
\newcommand{\OriginalPubYear}{1922}   % original publication date (if reprint)
\newcommand{\ThisEditionYear}{2027}   % current edition date

% end metadata


\newcommand{\BookTitlePage}{%
\thispagestyle{empty}

\vspace*{2cm}

\begin{center}

{\Huge\bfseries \BookTitle \par}
\vspace{0.5em}

{\Large\itshape \BookSubtitle \par}

\vspace{2.5cm}

{\Large \BookAuthor \par}

\vfill

{\small
\PublisherCity \par
\PublisherName \par
\vspace{1em}
Originally published \OriginalPubYear \par
This edition \ThisEditionYear
}

\end{center}

\clearpage
}


\usepackage{fontspec}
\setmainfont{EB Garamond}

\usepackage[
  paperwidth=5.5in, paperheight=8.5in,
  top=0.75in, bottom=0.95in,
  inner=1.0in, outer=0.7in % opposing margins for binding
]{geometry}

\usepackage{setspace}
\setstretch{1.25}
\usepackage{microtype}

\usepackage{titlesec}
\titleformat{\chapter}[display]
  {\normalfont\huge\bfseries}
  {\thechapter.} % or {} to suppress completely
  {0pt}
  {\Huge}

\titlespacing*{\chapter}{0pt}{-0.4in}{18pt}

\usepackage{fancyhdr}
\setlength{\headheight}{14pt}
\pagestyle{fancy}
\fancyhf{}
%\fancyhead[LE,RO]{\thepage} % page numbers: Left Even, Right Odd
%\fancyhead[RE]{\leftmark}   % chapter title on even pages
%\fancyhead[LO]{\rightmark}  % section title on odd pages
\fancyhead[LE,RO]{} % page numbers: Left Even, Right Odd
\fancyhead[RE]{}   % chapter title on even pages
\fancyhead[LO]{}  % section title on odd pages
\renewcommand{\headrulewidth}{0pt} % no top line
\renewcommand{\contentsname}{My Custom Title}

\begin{document}

\frontmatter
\BookTitlePage

\clearpage
\thispagestyle{empty}
\vspace*{\fill}

\begin{center}
\small
\EditionNote
\end{center}

\clearpage

\mainmatter

\chapter{An Opening Word}
Many voices had been speaking of eternal life, before the days of
the Son of man.  Especially pronounced had been the teachings of
the Egyptians that there is another world.   In their Acadian hymns
the Chaldaeans had dimly foretold a future life.  The belief of the
Parsees, as expressed in their Zend-Avesta, had included a place of
darkness for the evil soul and a reward for the good in the realm
of light.  The Hindus had declared, in their Rig-Veda, their
beautiful conception of the immortality of the soul, and had
written of a future "imperishable world, where there is eternal
light and glory."  The Grecian and Roman mythologies had voiced
their hope of blessedness for the shades of the departed.
Everywhere serious men had been asking as to the experiences beyond
the grave.  It was as if the Eastern world had become a vast
parliament chamber, wherein the nations were proclaiming their
different doctrines as to a future life.

In the midst of these varying and uncertain voices, Christ spoke
his authoritative message.  There was no wavering in his tone.
What the Oriental philosophers were guessing, he revealed; what the
Hebrew prophets had foreshadowed in their holy writings, he
unfolded in full light.  The ancient Vedic hymns, the oracles of
Greece, the Egyptian \textit{Book of the Dead}, anticipating by two
thousand years the Hebrew exodus--all these are naught compared
with the words of that inspired Teacher who spoke in Palestine.

In addition, Christ was himself the vital evidence of the
resurrection which he taught.  Against the assaults of doubt his
unique teachings are buttressed forevermore by his own return from
the land of silence.  In a short week after his words to Martha at
Bethany he had become, through his own rare experience, the
resurrection and the life.  Not the dead Buddha, nor the departed
Zoroaster, nor the vanished  Pythagoras  ever came back through the
opened door of the sepulcher, wearing the grave clothes of those
who sleep.  Human fancy had never dreamed of such a rapturous
denouement for faiths other than Christianity.  The resurrection of
the Lord is the crowning narrative with which the Gospels close.
It is a risen Christ who repairs the wastage of human decay and
death.  A voice above all those from Ind or Persia or the Nile
speaks henceforth in Judaea and the world concerning immortality.
The superlative Easter argument is the risen Christ himself.

\chapter{A Roman Quest}
"If one might only have a guide to the truth."--\textit{Seneca}.


On Scopus, the high mountain north of Jerusalem, the Roman camp was
pitched, that last autumn in the ministry of Jesus of Nazareth.  A
few years further on, if the warriors of the Emperor Tiberius could
then have foreseen the future, Titus was to quarter his famous
legions on that vantage point; and from its elevation he was to
hurl himself as a resistless battering ram against the Holy City.
But, on this autumn day, when these chronicles begin, no blare of
trumpets was summoning the Roman soldiery to arms; only the feet of
the camp sentinels, as they walked their appointed rounds, broke
the quiet of the sunlit afternoon.

That lithesome, cultivated, serious-minded young knight, Quintus
Cornelius Benignus, is standing on the height which overlooks the
great metropolis.  He is the son of Marcus Cornelius Magnus, that
Roman noble who is the intimate associate of the reigning Caesar,
and who has been a luxurious resident on the Palatine Hill since
his distinguished proconsulship in Africa.

     *     *     *     *     *

NOTE.--It is not from any time-marked Hebrew roll that this story
of Quintus is now taken.  He was of Roman blood, and his record is,
rather, to be found in the Latin literature of his time.  Well it
is when some new leaf is discovered among the musty folios,
reciting the saintly character and the triumphs of those who lived
when Christianity was new.  This record shows the worth of
consecrated life and service in the days when the luxurious Roman
state most needed a Christian citizenship.  But the lesson is none
the less for these last days, when the hope of the world is in the
creed of Quintus.

     *     *     *     *     *

By the side of Quintus is his fellow soldier Aulus.  They had spent
their boyhood together among the scenes of Rome; now they are
companions still, on this last Roman expedition to the district of
Judaea.  While the common soldiery are throwing their  dice  in
the  camp thoroughfare, these are speaking of more serious things.
The picture on which they look from lofty Scopus includes the
shining roofs of Jerusalem, the wooded Mount of Olives, and the far
landscape to the south and west; its undulations and brilliant
colorings no Roman artist might put upon the canvas.

With the autumn haze covering the extended panorama, Quintus says
first to his comrade:

"What the fates have in store for me, here in the city of
Hierosolyma, I am much wondering.  The day before our trireme
sailed from Brundisium for Tyrus I made a visit to the augur's
tent.  His prediction was that my journey hither would be followed
by strange consequences.  The flight of the birds through the air
did not reveal to him just what was to occur; but that something
eventful was to take place he was very sure.  What is to be my
fortune?"

"Your lot it may be," answers Aulus, "to perform some daring deed,
here in our Jewish campaign; and on your return to Rome you may
receive a great reward from the hand of Tiberius."

"In my mind this has been," replies Quintus; "before I left Rome I
had an audience with our divine Caesar, and he was pleased to say
that my fidelity here might bring me special recompense.  Yet would
that be satisfying?  I have seen the triumphal processions in the
streets of Rome, when heroes have been acclaimed; I have heard our
statesmen in the Senate hall, and prize the joys of oratory; I have
been served all my days by slaves in my father's palace, and know
the sweetness of the Falernian wine in the banquet room.  A
proconsulate, if I might come to that dignity, would be a high
honor to write in my life story.  But, my dear Aulus, would there
be content in this?  My restless soul seems crying out for some
better gift from the gods."

"It cannot be," continues Aulus.  "that your heart's love is
involved.  When our military movements bring the Roman knights to
Palaestina, in their pride of birth they do not wed the black-eyed
daughters of the Jews.  On your earlier expedition to Egypt you met
a princess of the land, but were not let to espouse that swarthy
maiden of the Nile.  The reward of love cannot be the experience of
which the augur spoke at Brundisium."

"Not so," says Quintus in response; "as I was leaving Rome, it was
the beautiful Lucretia who sent me forth with her rare farewell.
For my return from Palaestina she is now waiting; and under the
blue skies of Italia we are to wed.  I have been wondering,"
Quintus adds further, "if the augur, watching the flight of birds
there at Brundisium.  could mean that I am to fall by death, here
in Palaestina.  We have not come for battle, but to guard the
peace.  Yet it is easy for Atropos, that cruel fate, to clip the
slender thread of life and send men on to die land of shades.  If
this was what the augur meant, no Roman in the days of Tiberius has
ever set forth upon a more serious adventure."

"You are given to melancholy, this autumn afternoon, my comrade
Quintus," the other says; "you are feeling that sadness which comes
to men when the Dryads move over the earth and touch the leaves
into crimson and gold and brown."

"Not so," answers Quintus; "but I am remembering that I have come
into a land where a strange Teacher is speaking to men of a future
life.  Yet are men to live again?  I have seen the marble tombs on
the Appia Via where the Scipios, the Metelli, and so many more of
our great Romans lie asleep.  Shall I soon follow them?  Is it an
endless slumber?  What is it that the new Rabbi from Nazareth
means, when in the city yonder he speaks of another life?"

"A fig for your weird autumn fancy," responds Aulus; "down to the
streets of Hierosolyma we will go, and among their novel sights we
will forget your serious meditations."


They walk that afternoon as sightseers through the crowded Jewish
emporium.  The shops remind them, with all their contrasts, of the
marts of Rome, for men always and everywhere have the trader's
passion.  In the narrow streets of Jerusalem they see the stir of
many activities.  The workman is hammering his brass; the shoemaker
shapes his sandals; the flax spinner is winding his thread; the
scribe sits on his mat, and is ready for his writing.  In the shops
they see costly merchandise for sale--silks and jewels, fine linens
and perfumes, delicious foods and drinks.  These have been imported
from far Arabia and India; they have been brought from distant
Persia and Media.  With all their variety, no taste, however
fitful, need go unsatisfied.

What a motley crowd is on the streets!  They hear the Aramaic
speech of Palestine, which Quintus has been taught by his Athenian
tutor, and their ears also catch the accents  of other  foreign
tongues.  They meet traders from western Zidon, sailors from Crete,
bearded Idumaeans from beyond Judaea, and scholars from far
Alexandria.  Magnificent Jerusalem it is!  Yet destined soon to
fall.  For the day draws near when the Roman Titus shall weep on
Scopus over its fading splendors and then shall smite it to the
dust.

One purchase only does Quintus make.  In a shop where Egyptian
wares are sold he says to Aulus:

"Look on this scarab, this sacred beetle, which has been shaped by
some workman down in Thebae on the Nile.  We may be sure that no
people believes more intensely in a future life.  What compliment
they pay this physical frame of men when they hold that embalmment
restores to the soul its former body!  After the judgment of
Osiris, if their lives be true, the worthy shall enjoy the
companionship of the great god forever.  No other people wears such
a visible emblem of their faith in another life.  I will buy this
scarab for an amulet against accident and evil."


But where had the workman gone who once had shaped that token of
immortality?  Whither had vanished his carver's skill?  Where had
disappeared his projects and his dreams?  Quintus is not thinking
of any proconsulship he may win, or even of the love light in the
eyes of Lucretia, as he climbs again the heights of Scopus.  Rather
he is meditating on the departed maker of scarabs--and on the
destiny of the soul.  For ages the philosophers have been
speculating about the future life.  Familiar is Quintus with the
views of Laelius and Seneca, among the Roman inquirers, and with
the teachings of the great Grecians who have spoken in classic
Athens.  But now the question leaps to the front.  Quintus is in
the city where Ayran travelers and Persian magi and Egyptian
priests are busy telling their theories of immortality.  He is in
the very streets, besides, where a sandaled Teacher from Nazareth
is declaring that the dead shall live again.  If but half is true
that this strange Man is reputed to have said, no priest of Jupiter
has ever uttered at Rome so luminous a word.  Can it be that
Quintus himself shall see this Christus and hear his message?  If
so, his will be in very truth a momentous quest.
\chapter*{In Solomon's Porch}
\textbf{"Give me new consolation, great and strong, of which I nave never
heard or read."}--\textit{Pliny}.

With increasing frequency Christ was now speaking his prophecies of
the life immortal.  In his earlier ministry he had been dwelling
upon the presence of the divine kingdom in the earth, the practical
conditions for membership therein, and the inclusion of Gentile as
well as Jew in the gracious provision.  Novel were his words.
Whoever had heard his discourse on the Mount or the parable of the
lost sheep was rich beyond the modern sons of men.  But now, in the
closing period of his stay with mortals, he was more frequently
foretelling the life to come.  Like a footworn traveler drawing
near the homeland, he was keenly anticipating his return to the
spirit world.  Those who listened to him heard majestic intimations
of a celestial country which eye had not beheld.  Nor is it to be
thought that the Gospels, in their restricted pages, have recorded
half his words concerning the heavenly land.

Now comes the opportunity for Quintus himself to hear this new
Teacher of the Jews.  A messenger from Pilate, sent on an errand to
the headquarters at Scopus, brings the tidings that Christ is in
Jerusalem as a visitor at the Feast of Dedication.  Favored are
those who hear through the years the world's commanding voices;
beyond estimate is the high privilege now granted Quintus.

"I will hasten in to Hierosolyma," he says to Aulus, who is
detained by camp duties; "I will hear him for myself; and I will
bring you back report as to this latest prophet of immortality."


With his soldier's cloak about him, in protection against the
winter's chill, Quintus is away to Jerusalem.  The national Feast
of Dedication attracts his notice.  A courteous Hebrew explains to
him that the joyful festival commemorates the cleansing of the
Temple after its profanation by Antiochus Epiphanes, two hundred
years before.  The procession of pious Jews, carrying their palm
branches and marching to the heights of Moriah, the chanting of the
great Hallel within the imposing fane, the ascription of praise to
Jehovah all impress the keen-eyed soldier.

The enthusiasm of it all!  Though of other blood, Quintus clearly
feels the thrill of patriotism that stirs the multitude about him;
and he understands in some measure their impatient waiting for the
coming prince who shall deliver Israel.


But is this all?  Instead it is only the beginning of the wonders
which the serious Quintus is to witness.  Forth he passes to the
eastern cloister of the Temple, known then among the Jews as
Solomon's Porch, in memory of their illustrious king.  The
bystanders tell Quintus that it is built of a fragment of the first
Temple which Nebuchadnezzar had left standing.  As the soldier
looks down the far-reaching aisle, he sees a quadruple row of white
Corinthian columns, one hundred and sixty in number, and extending
a length of many hundred feet.  The vista is most amazing.
Accustomed though he has been all his days to the magnificence of
the Roman architecture, he yields in willing admiration to the
splendors of the Solomonic porch.

Then--he sees the Christ!  Walking through that forest of massive
columns is the superlative Jew of his times, and of all times.  For
now--when the voices of that winter day are still, and Solomon's
Porch has vanished where stood those blessed feet--there is no
earthly measurement by which to estimate the Man whom Quintus saw.


Among the throng that surround him hostile Pharisees challenge him
to tell them plainly if he be the foretold Messiah.  With impatient
hearts they have waited long for their redemption.  Let him say if
their deliverer has now come.  Then shall they throw off the yoke
of the detested Roman rule and renew their ancient monarchy with
enlarging influence and increasing splendors.

Memorable words in answer does Quintus hear.  The Stranger  puts
aside  the thought of the Jewish struggle for an earthly throne,
and turns in his fancy to the quiet pastures where feed the flocks.
He is a guardian Shepherd; Israel and all the world besides are his
cherished sheep.  Those who are truly his shall hear his guiding
voice, and shall follow him.  They shall never perish.  From the
hand of the Shepherd no vandal shall steal his own away.  How the
words thrill!  Sometimes Quintus has seen in the Judaean pastures
the keeper with his flocks, and knows how unchanging is his
fidelity.  It is as if this watcher in his devotion is anticipating
the faithfulness of the greater Shepherd.  How entrancing is the
lesson to this seeking soldier from beyond the Adriatic!

Then does the Christ add another word more surprising than the
rest.  To men who are his sheep he makes a promise that compasses
the furthest limit of the eternities.  Of such he says: "Unto those
who follow me I will give the Life of the Ages.  Beyond the tomb
they are to live on forevermore."  Nor to the Jews alone, amid the
maze of those Corinthian columns, does the coming Shepherd speak.
The listening Roman soldier, wearing the armor of the empire on the
Tiber, comes within the circle of his promise.  Into the face of
Quintus he looks and benignly says: "There are other sheep not of
the Jewish pasture, to whom I shall give this unending life.  I
covet your great empire as my own.  O soldier of the Caesars,
follow after me!"


Back to the camp on Scopus the soldier goes, moved to his deepest
soul.  Impossible it seems to longer worship the Roman gods.  When
he has described to Aulus the Feast of Dedication, he repeats the
words he has heard in the Temple cloister, and says in deepest
seriousness:

"Most unearthly is the Man on whom I have looked to-day.  In his
speech a divine patience, kindness, and dignity combine.  As for
the words he spoke, I cannot tell their moving power.  The sayings
of our noblest Romans are feeble in the comparison.  Never have I
heard another speak as he has done about a future world.  Truly, an
unequaled Man is this new Teacher who is abroad in Judaea."

Sleep is of little consequence that night.  Is the word of the
augur at Brundisium beginning to be fulfilled?  In his tent Quintus
is wondering through the long hours if, among his people on the
Tiber, the Shepherd shall not find some sheep to whom he will give
the unending life.
\c"He appeared to them alive again the third day, as the divine
prophets had foretold."--\textit{Josephus}.


How often have men missed the sight of great historic occurrences,
in their attention to the routine of life!  So it was that Quintus
did not witness the tragic events of that Passover week on which
human destiny was to turn.  To Tyre on the Great Sea he had gone,
to arrange for the landing of a new quota of troops from
Brundisium.  The commander at Scopus had chosen him for the
responsible mission, in token of his especial fitness.  The
compliment was pleasing.  But in his absence he was ever thinking
of the promise made by the Teacher in Solomon's Porch, that the
sheep who followed him should have eternal life.

Astir was all Jerusalem, when the knight returned to Scopus.  It
was on the morning after the Lord's resurrection.  That Roman
centurion who had been at Calvary reviewed for Quintus the fateful
happenings.  With pomp reminding of a Roman triumph the Christ had
entered David's city; after four days Iscariot had betrayed him
with a kiss; for blasphemy Pilatus, the procurator, had sentenced
him to the cross; they had put on him a scarlet robe in mockery;
they had hung him between two robbers on the hill of Golgotha; a
brutal soldier now at Scopus had won by lot his seamless robe, and
was jauntily displaying it as a trophy; an uncanny darkness had
covered the Judaean sky; the soldier Longinus had pierced the
sufferer's side; they had buried the dead Christ in the garden tomb
of the Arimathaean Joseph.  Monumental events were these--all new
to Quintus, but destined to be written indelibly in the calendars
of Christendom.

"More than this," continues the centurion, "an amazing rumor is now
abroad in the city that yesterday the dead Christus awoke from his
sleep and has been five times seen by his amazed disciples.  When I
beheld him yield up the ghost, I hailed his death as that of a
devout man, but little did I think that he was a God and would
return from the tomb.  The report says he has now come back.  On
swift wing the rumor has flown through Jerusalem and even into
Pilate's palace."


Down from the heights of Scopus the hurrying feet of Quintus carry
him to Jerusalem.  Doubts and wonderings and half-beliefs fill his
mind.  What if by any shadow of possibility the prediction of the
strange Teacher has been fulfilled, that he should return from the
dead on the third day?  Finding his way to Joseph's garden, Quintus
stands by an empty sepulcher.  There is a group of wondering
visitors near, and among them is one whose inviting face leads
Quintus to accost him.  Not frightened by the sword and armor of
the Roman knight, but assured by his candid look, the other answers
in the Aramaic which both can speak:

"Johannes is my name.  Till three years ago I was a fisherman, up
on the waters of Gennesaret.  Since then I have been a disciple of
this Man from Galilee.  In his company I have heard surprising
words and have felt a heavenly influence.  He was no ordinary
Teacher.  He was indeed from above."

"Is it true," asks Quintus in breathless words, "that your Master
has risen from the grave?  I have been away in Tyrus.  Now in the
Roman camp on Scopus I have heard that he has come forth from the
sepulcher.  What means such a marvelous report?"

"Yes, it is all true," John answers with his face aglow; "this is
the very sepulcher where our Lord was laid.  Your own sentries kept
guard before the tomb securely sealed.  But on the morning of
yesterday there was a shaking of the earth; some angelic visitants
rolled away the stone door of the grave; and our immortal Christus
came forth again.

"Astounding," Quintus interrupts in a whirl of words; "but did he
make any promise of another life for men, before he was put to
death?"

"He truly did," replies the disciple; "when we had eaten the
Passover supper with him, he spoke a word more marvelous than any
of your Roman teachers has ever uttered.  Into the spirit world he
said he was departing, to make ready a room in the Father's ample
house for those who were his own; and on his return he would take
them to be with himself.  Ever since our sad-hearted band have been
comforting themselves with this last promise in the upper room."

"None of our Roman gods has ever promised such a future." responds
Quintus; "but is this all?"

"No," answers the disciple; "on his cross our Christus spoke again
about another experience for men.  By his side was Dysmas, the
crucified robber, grieving for his faults and asking comfort.  When
the cross pain and thirst were over, our Lord replied, the outlaw
should walk with him among the bowers of the beautiful Paradise
beyond this world's horizon.  It was enough.  In this consolation
the tortured Dysmas passed on, with a smile of peace upon his face."

"Have you more wonders to tell?" presses Quintus, in his eagerness,
while the story of the cross begins to compel his judgment and call
for his heart's surrender.

Then, the consummation!  In ecstatic words John tells of the one
final and overmastering proof, in the thought of the eleven
disciples;

"Greatest of all, we have ourselves seen our Friend again.  Five
times already has he showed himself.  First, Mary of Magdala saw
him under the trees of the garden, and spoke with him; then the
other women met him and fell at his feet; next our fellow disciple
Petros saw him; then two of our band walked with him to outlying
Emmaus, and knew him as he broke bread at the journey's end; and
then last evening, he came to ten of us in the Passover room and
spoke his peace on us.

"Perhaps you have all seen a spectral form which has no real
existence," remonstrates Quintus, while all the time he is yielding
himself to the compelling story.

"It cannot be," responds the convincing John; "there have been too
many witnesses for that.  We have seen the very wound made by the
spear of Longinus; we have heard his familiar voice; we have
received his blessing.  Our number is our evidence; it cannot be
possible that all of us have been deceived.  It is surely he, O
Roman soldier, unless the senses of the women and of ten honest men
are far astray.  No other teacher of the East has ever come back
from the sepulcher.  Look and see for yourself.  Yonder is Joseph's
empty tomb.  The Christus is himself the evidence."


What can Quintus do, in the face of such proof as this?  He returns
to Scopus in wildest tumult.  Little does he say to Aulus, his
chosen friend.  The company of Longinus or the centurion he does
not seek.  The time has come--as it comes to all--when he must
commune with himself, and make the decision confronting every soul
that has heard the resurrection story.
\chapter{Cicero or Christ?}
"The name of Jesus can still remove distractions from the minds of
men."--\textit{Origen}.


Shall men believe in a future life because of Christ's return from
the grave?  Is his established resurrection at Jerusalem the
climacteric proof for immortality?  The problem is inescapable.
Every man is himself a judge; before every man the accumulated
evidence passes; for every man it is doomsday when he stands at the
point of decision.

In his sore perplexity Quintus says to himself that night, when he
has returned from his interview with the disciple John: "My soul is
like a traveler who halts at the point where two roads meet.  Great
issues depend upon his choice.  But while he hesitates may the
immortals, who watch over the destinies of men, guide his feet
aright."

Clearly defined are the alternatives before the Roman soldier.  On
the one hand are his ancestral beliefs, long established and deeply
cherished by the nation.  Nor does any man quickly toss aside the
faith of his fathers.  If belief is waning in the primitive
mythologies, and if the social life of the Empire is moved by
unrest and despair, the problem is to find a greater satisfaction.
There have been spoken many beautiful words by the Roman scholars
which are sweet premonitions of immortality.  Does not Quintus
remember that Cicero likens to heaven a port prepared, and prays
that he may sail thither with full-spread sails?  And if the gifted
Cicero has just gone tragically out of life, let it be hoped that
he has reached the harbor.

But on the other hand are the challenging and captivating words of
Christ.  Had he only spoken of the future life as an enthusiastic
Teacher, and then had passed to the perpetual slumber of the grave
like other philosophers of the time, he would be remembered long.
But, when he had spoken his words concerning immortality, he had
added, "I myself shall surely come back again."  From the evidence
which Quintus has heard in Jerusalem he has now fulfilled his
prediction.  He has put to scorn the fidelity of the Roman
sentinels at the tomb of Joseph; he has reversed the laws of
nature; he has appeared again, in unique proof that there is to be
a resurrection of the dead.  Wide is the difference between Cicero
and the Christ.  The one has spoken a mere opinion, so beautiful in
its phrase that it shall pass down into the future literature of
men.  The other has spoken a revelation, and then has returned to
prove that revelation true.  Which shall it be--Cicero or the
Christ?

But to accept the Jewish Teacher means earthly loss.  As he keeps
guard with himself through the night hours Quintus is wondering if
he shall incur the hostility of his father Marcus and shall be
forced to sacrifice his estates on the Palatine.  He fancies also
the grief of the fair Lucretia when she learns that he has chosen
an alien faith.  And he remembers, further, that in the choice of
the Christus he is joining a company on whom the Eastern world is
already casting its withering contempt.  Cicero or the Christus.
Which shall it be?


There are no struggles like the night wrestlings of the soul in
matters of religion.  What words can measure the divers arguments,
the opposing considerations, the conflicting emotions that shape
human choice?  Quintus stands at the point where soon--in the
progress of the new faith--Saul from Tarsus, Clement of Rome, and
so many more of the great spirits of that first era are to stand.
The wrestlings of the night!  Then foul demons are abroad; and then
God's good angels are descending the ladders of the sky.

Soon comes a great moment.  While the soul of Quintus is in wild
commotion, there falls upon him a mighty force which is not of
earth.  Coming he knows not whence, but not invading the department
of his will, it impels him to the Christ.  Transformed is this
Roman knight, who has been taught the doctrines of the Latin cult,
and whose nation can only feel disdain for a Galilaean who proposes
to revolutionize the ages.  The words of the augur at Brundisium
are having in truth a strange fulfillment.

As if the Man were present on whom he had looked in the Porch of
Solomon.  Quintus speaks his choice for the long eternities:

"Happen what may, I take thee, O Christus, for my Lord and Master.
I sacrifice my Roman knighthood for thee, if it shall be required.
I choose thee, because thou hast risen from the dead and hast
proven that there is another life for men."

Not Cicero, but Christ!  The Roman knight has made the great
decision.
\chapter*{The Vision of the Risen Christ}
\textbf{"After that, he was seen of above five hundred brethren at
once."}--\textit{Paul}.


Once for himself was Quintus to see the Lord, before his departure
heavenward.  When midnight hours afterward came to him in Italy,
the memory of that vision was golden.  When, among the temples of
the gods in pagan Rome, men challenged his belief, his sufficient
answer was: "With mine own eyes I have seen the risen Teacher who
has revealed immortality to men."  So did the first disciples of
the faith who bore its weightiest burdens, enjoy its highest
privilege.

It was the disciple John who told Quintus of the opportunity to see
the risen Lord.  In an hour of fellowship at Jerusalem--when the
knight had confessed his new allegiance--John spoke of the Master's
wish.  The disciples who were in the city and its environs were to
gather in Galilee with those from that upper district.  Once more
would their Lord show himself to all who believed on him, and would
speak with them.  Nor did Quintus ever cease to rejoice that he was
reckoned worthy to look that day on the Conqueror of death.

With light feet the Jerusalem company, some six score in number,
made the journey north to Galilee.  One subject only was on their
lips, as they followed the road through Samaria to Kurn Hattin,
near the Sea of Tiberias.  Here the Lord at the opening of his
mission had spoken his nine blessings to needy mortals; most
fitting it now was that on this memorable hillside he should utter
his farewell to those who had come to believe on him.  Thus would
the circle of his teachings end where it had begun.  Bright was the
picture.  The glint of the sunlight on the Galilaean sea so near at
hand, with the uncounted flowers of the spring-time that covered
the lower plains, lent a charm to the scene that Quintus remembered
always.


At the outset the Roman convert is impressed with the goodly number
of those first disciples.  They are not twelve or six score, but
many more.  They greet each other with the salutation, "Peace be to
you," and then they rapturously add, "To-day we shall see our
Lord."  In that intimacy which should always mark the followers of
Christ, they give Quintus their welcome; and at once he feels
himself among a congenial brotherhood.

One is by name Nicodemus, a member of the Great Sanhedrin.  Another
is one Bartimaeus, from southern Jericho, whose finger tips have
been his eyes, till the Lord has healed his blindness.  A third has
been a demoniac among the hills of the Gergesenes, and has been a
wandering and truculent challenge to his times.  A woman is there
from Jacob's well, with Salome and Susanna and the virgin mother
herself.  They are from southern Bethlehem; they have come from the
wild hills of Peraea, beyond the Jordan; many are from Galilee,
where Christ has found so many devoted followers.  All these, as
well as the immortal eleven who have composed the inner circle of
the Master's associates.

Two other peculiar disciples does Quintus see, both of whom have
been raised from the dead.  Lazarus has come, who has so often
welcomed the Lord to his home in Bethany; and with him are the
sisters, of whom one has heard the Teacher say.  "Whosoever liveth
and believeth in me shall never die."  The other is a young
vineyard keeper from the neighboring village of Nain, whom Christ
has restored.  His word to Quintus is:

"Last year I sickened with a fever and passed through the door of
death.  They were carrying me out for burial, and my widowed mother
was weeping as one weeps who has lost her only son.  The Master
halted the mourners, and called me back to earth.  I have never
told of the wonders which I saw in the spirit world; it would not
be lawful.  But I have been in the great spaces beyond the stars,
and know that the tomb is only a resting place for a little sleep."

"How many disciples are there here?" Quintus asks of the good John.
To which question the other answers:

"Over a half thousand.  It has been our Master's wish that every
disciple of his throughout the land should come to this meeting
place.  Unto all he would show himself once more, before he returns
to the upper life.  So they shall have a glad memory of his face,
and shall be strengthened in their coming tribulations by the hope
of immortality."

Then suddenly--the risen Lord has come!  The marvel of it!  The
splendor of it!  While the five hundred are talking together, the
air grows luminous with his presence.  Out of the invisible he
appears.  As suddenly he comes as Aurora in her chariot drives up
the eastern sky and brings in the shining day.  When the company
have fallen on their faces and have adored their Master, in the
hush that follows he gives them a great commission:

"You are to go forth." he says, "and herald my gospel to the world.
Let there be no laggards in your company.  It is a lifelong charge.
There is a task for Petrus and Johannes, for Philippus and
Mattheus, and for all.  You are to look for disciples everywhere.
You are to proclaim the message of repentance.  You are to give
them the waters of baptism, in the name of the God triune.  You are
to declare to sad-hearted men the promise of eternal life, until I
shall come again to take men to myself."

That honorable commission!  It was in coming days to stir the souls
of apostles and quicken the feet of missioners and fire with zeal
earth's coming reformers.  Nor does Quintus forget that he too has
his charge.   In the city on the Tiber is to be his task.  To his
home circle, to priests in the temples of the gods, and even to the
royal Tiberius he is to herald the gospel of the resurrection.  His
vision of the risen Lord is the measure of his opportunity.

Then the Master looks into his very face, and remembers him as the
Roman knight he had seen in the Porch of Solomon.  The half
thousand disciples on Kurn Hattin prostrate themselves to the
earth; and in their acclaim the soldier joins his voice, "Rabboni!
Rabboni!  Our great Master!"  Then departs the Christ, and back to
their homes they go, evermore to comfort themselves with the vision
of their risen Lord.


Soon afterward their Rabboni goes from earth.  Out beyond the hill
of Olivet he walks one day with his eleven.  In their last words
together he reminds them again that they are to be his heralds to
the eastern world.  A cloud gathers above their heads, like some
halting chariot, and he is gone forever from human sight.  Yet only
in the distance it seems a cloud.  For John afterward says to
Quintus that it was in reality a phalanx of ten thousand angels,
robed in whiteness and sent to convoy the Son of God to glory
everlasting.

With Quintus that visit to Kurn Hattin shaped all his future.  His
Master's countenance had seemed to him more wonderful than any face
which the gifted Phidias had ever carved in stone.  But never in
after days could he worthily tell to Lucretia the vision he had
seen.  Only in one poor sentence could he sum it up: "I have seen
for myself the risen and ascending Lord."
\chapter{Christ's Witnesses at Rome}
"A great multitude."--\textit{Tacitus}.


With jubilation Quintus sees again the shores of Italy rise over
the Adriatic, and finds himself once more in his beloved Rome.  The
center of magnificence and power it seems.  Alter clamorous public
greetings in the Forum, there comes another welcome which happens
only in a returning soldier's life.  In the palace of Marcus the
kindred of Quintus are gathered, and Lucretia also is in the
circle, to hear his great adventure.

"How wonderful it seems," the knight begins: "so many times have
your faces come to me in my dreams, but now I am fully awake and
see them once again.  Hail to you all!  When I was sailing away
from Brundisium, the augur foretold for me an unusual experience.
In the Jewish life beyond the Sea I have learned much, if that were
the fulfillment.  But, most of all, I have come back with a new
religious faith.  In Judaea, as you must have heard, a certain
Galilaean has called himself the Son of the one true God.  He has
spoken of a future life for men; and he has now risen from the
grave, after his torture on a cross, to prove his doctrine true.  I
now believe in him, as the interpreter of the future life.
Forevermore he is my High Priest, and not the great pontifex in the
temple of your Jupiter."

Brave words they are.  There in the great hall of Marcus, with the
sunlight shining on the gorgeous palaces of the Caesars, the Temple
of Apollo, and all else which crowns the Palatine, the noble
Quintus confesses his new belief.  Come what will the consequences!

Then, while they hear in amazement, he further says; "Most inviting
is this new creed.  Our wise Roman scholars, as well as those in
Greece, have only been guessers about the future life.  But the
Christus speaks as one who has come from the heavens.  Those who
keep his commandments are to dwell with him forevermore in eternal
joy.  Everywhere through Judaea men are becoming his followers, and
the wide world is to believe on him.  Perhaps you also, my
cherished ones, will come to accept his teaching of the future
life."

So Quintus speaks, with his vibrant voice and with a strange light
on his face.  Wonderingly they hear the tidings that he brings--the
recital of the greatest happening that can ever befall a man.  Not
deriding their valiant soldier, and not withholding their wealth of
love from one who has come safely back to them, they watch the
changes in his life.

"I do not care," he says, "to loiter in the baths of Agrippa and to
hear from the idlers there the gossip of the hour.  The
gladiatorial struggles in the Circus Maximus and the comedies in
the theaters have lost for me their relish.  For the civic rewards
which Tiberius gives his favored ones I have no wish.  Senatorships
and proconsulships are like the dust in the apothecaries' scales.
I have seen the risen Lord!"

Influential is such a life on the home group of Quintus.  With his
pride of birth and his great properties, Marcus becomes a believer.
A conversion it is which is the surprise of Rome.  The rare
Lucretia, as well, receives the truth.  At times, before she has
called herself a disciple, Quintus escorts her to the worship of
the Roman Christians.  Their captivating speech, their holy love
for one another, their rapturous faces move her deepest heart.
Till, one day, when Quintus has been telling her of the womanhood
in Judaea which the Christ has ennobled, she replies:

"I believe it all, O Quintus.  Of late into my heart an untold
peace has come.  All things are changed for me.  The sunlight is on
the hills!"  It is her open confession.  Lucretia is thenceforth
enrolled among the Roman saints of whom the world was not worthy,
and who looked for the life to come.


In the fellowship of the Roman church--already founded and rapidly
enlarging--Quintus finds his pleasure.  A few are Jews from the
ghetto beyond the Tiber, till the persecution of Claudius drives
them forth.  More are of the varied nationalities met in that
commercial and luxurious center.  Most are of plebeian blood.
There are smiths and mechanics; there are stone cutters, workers in
mosaics, and decorators.  There are slaves from the very palace of
Tiberius.  There is Amon from Egypt, who sells his jewelry down in
the Nova Via.  There is Polemon, the Grecian shopkeeper, in the
Clivus Victoriae.  There is Onesimus, the servant of Philemon, from
Colossae.  There are Amplias and Epaenetus and Stachys, the
particular friends of the Gentile apostle.  There is, as well,
Pomponia Graecina, that woman of noble blood, who accepts the
Christ.  An ever-increasing company it is.

In their assemblies, on the first day of the week, Quintus has his
influential place.  He listens to the reading of the older
Scriptures; he celebrates with the gathered company the eucharistic
suppers and agapae; he keeps with them the Easter celebration, in
memory of Him who shall give them eternal life.  In emblem of their
faith the sign of the fish is on their evening lamps.  Theirs is a
sterling citizenship.  The wanton metropolis of the Caesars is
blessed immeasurably by the company of these who follow the risen
Lord.

It is after the midcentury that the great Paulus, having met with
shipwreck on Melita, draws near to Rome.  Quintus leads the company
that goes out southward forty miles, to welcome the Christian
traveler.   At Appii Forum, that common town with its bargemen and
its tavern keepers, they give the kiss of welcome to a little bent
and gray-haired Jew, who shall go down into history as Christ's
most illustrious apostle.  The faithful Luke is his companion.
Along the famous highway of the Via Appia, where emperors and
warriors, scholars and Oriental tradesmen have walked, Quintus
escorts their guest.  Past the tombs of the Roman great, by
uncounted statues, past suburban villas they go, until, through the
Porta Appia, the holy prisoner, chained to a Roman guard, finds
himself in the city of the Caesars.

One rare privilege the Roman knight then envoys.  In his hired
house, near the Pretorian camp, Paul speaks without interruption
his words of grace.  The doctrines he had before written to the
Roman church he now explains; the wish he had made to see them face
to face now expresses itself in words of love.  The flood tides of
his eloquence move resistlessly on, as he interprets the new faith
and speaks of Him who is to give them eternal life.  Quintus is
enriched by his frequent association with the peerless soul.  Nor
did he have a prouder thing to say, in the days to come, than to
declare, "I heard great Paulus tell of the life immortal."

But how fares our knight when persecution comes?  Through the years
he has been bravely declaring the Christian doctrine of the eternal
life to priests in the temples, to Roman nobles, to all most
hostile.  But his wealth and social standing, as well as the
emperor's favor, now insure his safety.  His father Marcus has long
since passed on, in hope of the heavenly life.  Having wedded the
graceful Lucretia, when an apostle was in Rome to speak their
nuptials, he has her efficient counsel in the testing times.

"Look! look! Lucretia," he cries, one evening; "through the lower
city the flames are running like unbridled horses.  There is danger
that all Rome may go to ashes."

For nine long days they watch the sweep of the lurid flames.  The
light shines out like a signal torch, to mark an emperor's folly.
Then the undeserved charge that they have lit the flames brings on
the martyrdom of the Roman Christians.  Sometimes Quintus and
Lucretia are able to soften the trials of the sufferers, by
permission of the capricious Nero.  To old Chilo, the Grecian,
before he meets his doom, they unfold the promise of eternal reward
in the Father's house.  The hope of immortality they carry to those
who go to the lions, at the emperor's whimsical command.  And the
glorious company of martyrs passes singing to the skies, because of
their consoling words.

Down into the dungeon of the Mamertine they are permitted once to
go, to visit Paulus.  But he needs not their consolation.  Rather
he is the comforter.  With the poise of a conqueror he bids them
not to mourn for him: he is going to the Lord in the unending life.
Over their bowed heads he stretches his aged hands, in apostolic
benediction.  Soon ends his imprisonment.  At \textit{Tre Fontane}, in a
few days more, his weary body rests; but his immortal spirit mounts
beyond the stars.

At last the Christian knight comes to the crossing.  The prediction
of the augur at Brundisium has been strikingly fulfilled.  Matured
in all the graces, he is like the ripened Chian clusters that await
the vintager in the autumn days.  The friends of Quintus have gone
before; as the old century wanes, the old man is to follow them.

"My time has come to go," he says one day; "the portals of eternal
life and joy I see swinging open wide.  I shall pass through the
gates, because my ascended Lord has gone in before me to prepare my
dwelling place.  With him as my Teacher I believe in the life
immortal."


In the Roman catacombs, those most remarkable testimonies to the
eternal life, his resting place may be found.  The sign of the fish
is on his stone.  Its time-eaten inscription is still legible,
among the many which tell of the early Christian expectation and of
all future Christian hope:


"HERE RESTS THE DUST OF QUINTUS, OF NOBLE BLOOD; IN THE FAITH OF
THE ASCENDED LORD HE HAS ENTERED UPON THE ETERNAL LIFE."

\backmatter
\chapter*{Afterword}

\indent Easter is, above all, a celebration of hope. The story of Christ does not end upon the Cross or in the silence of a sealed tomb. On the third day the stone was rolled away, the tomb was found empty, and sorrow gave way to the words that have echoed through the centuries: He is risen. An Easter Disciple, first published in 1922, seemed a fitting little work to reprint for the season. Through the story of Quintus, Arthur Benton Sanford looks upon the events surrounding the Crucifixion and Resurrection through the eyes of one who must encounter them for himself. Though his story is fiction, the truth at its heart is the foundation of Easter: Christ was crucified, Christ died, and Christ rose again.\\

\indent This edition was reprinted and hand-bound as part of my bookbinding hobby and as a small Easter gift for friends and family. There is a quiet pleasure in taking an old text, setting it once more upon the printed page, folding and sewing its gatherings, and giving new physical form to words written more than a century ago. I offer this little volume in that same spirit—as a remembrance of Easter and of the hope proclaimed by the empty tomb. The paper will age, the thread may loosen, and the binding will bear the marks of the hands that hold it; but the message it carries remains: \textbf{He is risen indeed.}

\vspace{1.25em}

\noindent
--- Arthur M. Strutzenberg\\
Easter, 2027

\end{document}